% Evans 10.3: printed pp.554--566; English PDF pp.571--583.
% Chinese section start: printed p.391 / scan p.409.
\section[Control Theory, Dynamic Programming]{Control Theory,\texorpdfstring{\\}{ }Dynamic Programming}
\label{sec:10-3}
\subsection{Measurable controls and deterministic responses}
Let $A\subset\R^m$ be nonempty and compact. Suppose
$f:\R^n\times A\to\R^n$ and $r:\R^n\times A\to\R$ are continuous,
and uniformly in $a\in A$ satisfy
\[
 \begin{aligned}
 |f(x,a)|&\leq M,& |f(x,a)-f(y,a)|&\leq L_f|x-y|,\\
 |r(x,a)|&\leq R,& |r(x,a)-r(y,a)|&\leq L_r|x-y|.
 \end{aligned}
\]
Let $g$ be bounded and $G$-Lipschitz. Fix a horizon $T>0$.
An admissible control $\alpha:[t,T]\to A$ is measurable. Its state is
an absolutely continuous solution of
\begin{equation}\label{eq:10-control-ode}
 X'(s)=f(X(s),\alpha(s))\text{ a.e.},\qquad X(t)=x.
\end{equation}
\Needspace{50mm}
\begin{lemma}[Response estimates; details for source \S10.3.1]
Every admissible control has a unique response, and
\[
 |X(s)-x|\leq M(s-t),\qquad
 |X^{x,\alpha}(s)-X^{y,\alpha}(s)|\leq e^{L_f(s-t)}|x-y|.
\]
\end{lemma}
\begin{expandedproof}[the integral equation accepts measurable controls]
For fixed $\alpha$ use the map
$Z\mapsto x+\int_t^s f(Z(q),\alpha(q))dq$ on continuous paths.
The integrand is measurable and bounded; the map returns continuous
paths and is a contraction on intervals with $L_fh<1$.
Iterating these intervals gives an absolutely continuous solution on
$[t,T]$. Its displacement bound follows by integrating $|f|\leq M$.
For two responses with the same control, subtract the integral equations
and apply Gronwall. This gives the second estimate and uniqueness.
No differentiability in time of the control is required.
\end{expandedproof}
The cost and value are
\begin{equation}\label{eq:10-value}
 \begin{split}
 J_{x,t}(\alpha)&=\int_t^T r(X(s),\alpha(s))ds+g(X(T)),\\
 V(x,t)&=\inf_\alpha J_{x,t}(\alpha).
 \end{split}
\end{equation}
An infimum does not itself assert existence of an optimal control.
The following arguments use controls within an arbitrarily prescribed
positive error of the infimum.
\Needspace{42mm}
\PDEAnnotation{紧控制集保证什么}{连续性与紧性保证 Hamiltonian 的逐点最小值能取到。但这并不自动给出整个时间区间的最优控制，也不自动保证反馈场的闭环 ODE 可解；动态规划证明只使用近似最优控制。}

\Needspace{70mm}
\subsection{The dynamic programming principle}
\Needspace{55mm}
\begin{theorem}[DPP; source \S10.3.2, Theorem 1]
For $0<h\leq T-t$,
\begin{equation}\label{eq:10-dpp}
 V(x,t)=\inf_{\alpha|_{[t,t+h]}}
 \left\{\int_t^{t+h}r(X(s),\alpha(s))ds+V(X(t+h),t+h)\right\}.
\end{equation}
\end{theorem}
\begin{expandedproof}[restriction and concatenation give opposite bounds]
Call the right side $W$. Any full control has a prefix up to $t+h$
and a continuation thereafter. Its continuation cost is at least the value
at its reached state. Thus every full cost is at least $W$, giving $V\geq W$.

Conversely, fix any prefix, with reached state $z$. For $\delta>0$
choose a continuation whose cost is at most $V(z,t+h)+\delta$.
Concatenate the two measurable controls. Uniqueness of the state equation
makes its trajectory exactly the concatenation of their responses.
Therefore $V(x,t)$ is at most the displayed prefix cost plus $\delta$.
Take the infimum over prefixes and let $\delta\downarrow0$, obtaining
$V\leq W$. A continuation is chosen for a single reached state at a time;
no measurable selection depending simultaneously on every state is needed.
\end{expandedproof}
\PDEBodyDiagram{10.3}{A prefix and its continuation form one admissible control.}
{Restricting a full control gives a continuation cost at least $V(z,t+h)$. Conversely, for a fixed prefix reaching $z$, choose one $\delta$-optimal continuation and concatenate it. The two comparisons prove the equality in DPP. The curves are schematic state trajectories, not a claim of optimality for the drawn paths.}
{ch10-dpp}
\Needspace{50mm}
\begin{lemma}[Regularity of the value; source \S10.3.3 lemma]
The value is bounded by $RT+\|g\|_\infty$. Uniformly in $t$ it is
$K_T$-Lipschitz in $x$, where
\[
 K_T=Ge^{L_fT}+L_r\int_0^T e^{L_fs}ds.
\]
It is $(R+K_TM)$-Lipschitz in time and $V(x,T)=g(x)$.
\end{lemma}
\begin{expandedproof}[compare identical controls before taking infima]
For the same control started at $x,y$, the response estimate bounds the
cost difference by
\[
 \left(G e^{L_f(T-t)}+L_r\int_0^{T-t}e^{L_fs}ds\right)|x-y|.
\]
This holds for every control, so taking infima yields the same bound
for $|V(x,t)-V(y,t)|$. In DPP every reached state satisfies
$|X(t+h)-x|\leq Mh$ and every running integral has absolute value
at most $Rh$. The expression inside its infimum is therefore between
$V(x,t+h)-(R+K_TM)h$ and $V(x,t+h)+(R+K_TM)h$.
Taking its infimum gives the time estimate without assuming a smooth
value function or an optimal control. At $t=T$ the running interval is
empty and $X(T)=x$, so every control gives cost $g(x)$.
\end{expandedproof}
\Needspace{42mm}
\PDEAnnotation{先用同一控制比较}{估计两个初值的价值差时，对它们使用相同控制，先用 Gronwall 控制轨道差，再对代价取下确界。直接把两个未知最优控制相减没有这样的稳定性依据。}

\Needspace{70mm}
\subsection{The terminal HJB equation and its time orientation}
Define the control Hamiltonian
\begin{equation}\label{eq:10-control-H}
 H_c(p,x)=\min_{a\in A}\{f(x,a)\cdot p+r(x,a)\}.
\end{equation}
It is concave in $p$, since it is a minimum of affine functions, and
\[
 \begin{split}
 |H_c(p,x)-H_c(q,x)|&\leq M|p-q|,\\
 |H_c(p,x)-H_c(p,y)|&\leq(L_f|p|+L_r)|x-y|.
 \end{split}
\]
For a terminal problem written as $V_t+H_c(DV,x)=0$, the source
uses the following convention: an upper contact has residual
$\phi_t+H_c(D\phi,x)\geq0$, and a lower contact has residual at most zero.
Equivalently this is the usual viscosity convention applied to
$-V_t-H_c(DV,x)=0$.
\Needspace{50mm}
\begin{lemma}[Time reversal; source remark after Theorem 2]
Set $w(x,\tau)=V(x,T-\tau)$. Then the terminal convention for $V$
is exactly the initial-value convention for
\[
 w_\tau+K(Dw,x)=0,\qquad K(p,x)=-H_c(p,x),\qquad w(x,0)=g(x).
\]
In particular $K$ is convex in $p$ and satisfies \eqref{eq:10-H-lip}.
\end{lemma}
\begin{expandedproof}[transform the test as well as the equation]
If $w-\psi$ has an upper contact at $(x_0,\tau_0)$, then
$V-\phi$, $\phi(x,t)=\psi(x,T-t)$, has an upper contact at
$(x_0,T-\tau_0)$. Now $\phi_t=-\psi_\tau$, $D\phi=D\psi$.
The terminal inequality $-\psi_\tau+H_c(D\psi,x)\geq0$ is precisely
$\psi_\tau-H_c(D\psi,x)\leq0$. Lower contacts give the other direction.
The implication is reversible. Negating the minimum of affine functions
makes a maximum of affine functions, proving convexity of $K$.
\end{expandedproof}
\PDEMarginDiagram{10.M}{时间方向与测试符号}
{图中残差均按方程所写的时间变量计算。令 $\tau=T-t$ 后，$\partial_t=-\partial_\tau$，Hamiltonian 同时变为 $-H_c$；不能只改变时间端点而保留所有符号。}
{ch10-time-margin}
\Needspace{55mm}
\begin{theorem}[Value function solves HJB; source Theorem 2]
The value in \eqref{eq:10-value} is the unique BUC solution of
$V_t+H_c(DV,x)=0$, $V(x,T)=g(x)$, with the terminal convention just
specified. Equivalently its time reversal solves the corresponding
initial-value problem with Hamiltonian $K=-H_c$.
\end{theorem}
\begin{expandedproof}[all constant controls versus an $h^2$-optimal prefix]
Normalize an upper test by $V(x_0,t_0)=\phi(x_0,t_0)$.
For any fixed $a\in A$, use the constant control on $[t_0,t_0+h]$.
The displacement bound keeps its path in the contact neighborhood when
$h$ is small, independently of $a$. DPP and upper touching give
\[
 \begin{split}
 0&\leq\int_{t_0}^{t_0+h}r(X(s),a)ds
              +V(X(t_0+h),t_0+h)-V(x_0,t_0)\\
 &\leq\int_{t_0}^{t_0+h}
       \bigl(\phi_t+f(X(s),a)\cdot D\phi+r(X(s),a)\bigr)ds.
 \end{split}
\]
Divide by $h$ and let $h\downarrow0$. The limit inequality holds for
every $a$, hence
$\phi_t(x_0,t_0)+H_c(D\phi(x_0,t_0),x_0)\geq0$.

At a normalized lower contact choose a prefix within $h^2$ of the
DPP infimum. The lower touching inequality now yields
\[
 h^2\geq\int_{t_0}^{t_0+h}
 \bigl(\phi_t+f(X(s),\alpha(s))\cdot D\phi+r(X(s),\alpha(s))\bigr)ds.
\]
Continuity in $(x,t,a)$ on a compact neighborhood times $A$ makes
replacement of $(X(s),s)$ by $(x_0,t_0)$ have a uniform $o(h)$ error.
The resulting integrand is bounded below by
$\phi_t(x_0,t_0)+H_c(D\phi(x_0,t_0),x_0)$ for every control value.
After division by $h$ and passage to the limit we obtain a residual
at most zero. The error must be $o(h)$: a fixed $\delta$-optimal control
would leave the uncontrolled quotient $\delta/h$.

The value is BUC and takes the terminal data by the regularity lemma.
Time reversal converts it to an initial-value solution for $K$, whose
structural estimates allow comparison in \S10.2. This proves uniqueness.
\end{expandedproof}
\Needspace{42mm}
\PDEAnnotation{近似最优误差要随步长缩小}{在下接触证明中取 $h^2$-最优前缀，再除以 $h$，误差才趋零。上接触使用每个常值控制；下接触使用近似最优控制，二者的量词不同。}

\Needspace{70mm}
\subsection{Verification and the limits of feedback synthesis}
\Needspace{55mm}
\begin{theorem}[Smooth verification; editorial completion of source discussion]
Suppose $v\in C^1(\R^n\times[0,T])$ satisfies
$v_t+H_c(Dv,x)=0$ and $v(x,T)=g(x)$. Then $v(x,t)\leq J_{x,t}(\alpha)$
for every admissible control. If an admissible $\alpha_*$ with response
$X_*$ attains the Hamiltonian minimum along its path a.e., then its cost
equals $v(x,t)$ and it is optimal; in particular $v=V$ at that starting point.
\end{theorem}
\begin{expandedproof}[integrate the nonnegative verification residual]
The chain rule for an absolutely continuous path gives
\[
 {d\over ds}v(X(s),s)+r(X(s),\alpha(s))
 =v_t+f(X(s),\alpha(s))\cdot Dv+r(X(s),\alpha(s))\geq0
\]
a.e. Integrate to $T$ and use the terminal value $g$ to obtain
$J_{x,t}(\alpha)\geq v(x,t)$. If the minimum is attained along $X_*$,
the residual vanishes a.e. and equality holds. All evaluations are along
a bounded path on a finite interval, where the derivatives of $v$ are bounded.
This verifies an admissible path; it does not establish solvability of a
closed-loop ODE for an arbitrary minimizing feedback field.
\end{expandedproof}
The pointwise prescription
$a_*(x,t)\in\operatorname*{argmin}_{a\in A}\{f(x,a)\cdot DV(x,t)+r(x,a)\}$
is conditional on the gradient and on an admissible resulting trajectory.
At nonsmooth points $DV$ may not exist. Even with a pointwise minimizer,
measurability and closed-loop existence require additional arguments.
\Needspace{42mm}
\PDEAnnotation{验证定理与控制构造}{验证定理说：若已有合适的最小化轨道，则它最优。它没有证明任意逐点选出的反馈都可积分成轨道；尤其在价值函数不可微处，直接写 $DV$ 没有定义。}

\Needspace{70mm}
\subsection{Hopf--Lax: compact minimizers and viscosity inequalities}
Return now to an initial-value Hamiltonian $H:\R^n\to\R$, finite,
convex and superlinear, with $g$ bounded and $G$-Lipschitz. Let
\[
 L(q)=\sup_p\{p\cdot q-H(p)\},\qquad
 Q_tg(x)=\min_y\left\{tL\left({x-y\over t}\right)+g(y)\right\}.
\]
Convex duality from \S3.3 gives finite continuous convex superlinear $L$
and $H(p)=\sup_q(p\cdot q-L(q))$. In particular all displayed minima
exist by coercivity. These assumptions include $H(p)=|p|^2/2$, which
is outside the global momentum Lipschitz hypothesis of \S10.2.
\Needspace{55mm}
\begin{lemma}[Semigroup, speed bound, and regularity; details for source Theorem 3]
One has $Q_hQ_tg=Q_{h+t}g$. Each $Q_tg$ is $G$-Lipschitz in space.
All minimizing velocities $q=(x-y)/t$ lie in a ball independent of $x,t$.
With $Q_0g=g$, the function $U(x,t)=Q_tg(x)$ is BUC on every finite slab
and is Lipschitz in time.
\end{lemma}
\begin{expandedproof}[Jensen and the competitor with zero velocity]
For a two-stage path from $y$ through $z$ to $x$, Jensen gives
\[
 tL\left({z-y\over t}\right)+hL\left({x-z\over h}\right)
 \geq(t+h)L\left({x-y\over t+h}\right).
\]
Equality is obtained by $z=(tx+hy)/(t+h)$, so infimizing proves the
semigroup identity. Translating $x$ and $y$ by the same increment
changes only $g(y)$, proving the spatial Lipschitz estimate.

At a minimizer compare with $y=x$. Since $g$ is $G$-Lipschitz,
$L(q)-G|q|\leq L(0)$. Superlinearity makes this sublevel bounded,
say $|q|\leq R_G$. The bound is uniform even as $t\downarrow0$.
For any $G$-Lipschitz $v$, set
$m_G=\inf_q(L(q)-G|q|)>-\infty$. Then
\[
 v(x)+hm_G\leq Q_hv(x)\leq v(x)+hL(0).
\]
Together with the semigroup this gives a uniform time Lipschitz constant
$\max\{|m_G|,|L(0)|\}$ and convergence to $g$ at $t=0$.
The same estimate with $v=g$ bounds $U$ on finite slabs.
\end{expandedproof}
\Needspace{55mm}
\begin{theorem}[Hopf--Lax viscosity solution; source \S10.3.4, Theorem 3]
Under the preceding assumptions $U(x,t)=Q_tg(x)$ is the unique BUC
viscosity solution of $U_t+H(DU)=0$, $U(x,0)=g(x)$.
\end{theorem}
\begin{expandedproof}[bounded minimizing velocities close the lower test]
At an upper contact at $(x,t)$, choose any fixed velocity $q$ and
$0<h<t$. The semigroup gives
$U(x,t)\leq hL(q)+U(x-hq,t-h)$.
Upper touching then gives
$\phi(x,t)-\phi(x-hq,t-h)\leq hL(q)$.
Divide by $h$ and let $h\downarrow0$ to get
$\phi_t+D\phi\cdot q-L(q)\leq0$ for every $q$.
Taking the supremum in $q$ yields $\phi_t+H(D\phi)\leq0$.

At a lower contact, let $q_h$ minimize the short-step semigroup formula,
so $U(x,t)=hL(q_h)+U(x-hq_h,t-h)$. The spatial Lipschitz estimate
for $U(\cdot,t-h)$ bounds $|q_h|\leq R_G$ by the preceding lemma.
Thus the earlier point stays in the contact neighborhood and Taylor's
first-order remainder is $o(h)$ uniformly in $q_h$. Lower touching gives
\[
 hL(q_h)\leq\phi(x,t)-\phi(x-hq_h,t-h)
 =h(\phi_t+D\phi\cdot q_h)+o(h).
\]
Since $D\phi\cdot q_h-L(q_h)\leq H(D\phi)$,
$\phi_t+H(D\phi)\geq-o(1)$, which proves the lower-test inequality.
Initial data and BUC regularity were proved in the lemma.

For uniqueness compare $U$ with any other BUC solution in both directions.
The extension at the end of \S10.2 applies because $U$ is uniformly
spatially Lipschitz and $H$ is continuous and independent of $x$.
This explicitly supplies the comparison step without assuming a global
Lipschitz bound for the superlinear $H$.
\end{expandedproof}
\PDEBodyDiagram{10.4}{Two minimizing origins can create a viscosity corner.}
{Take $H(p)=p^2/2$, $t=1/2$, and $g(y)=\min\{|y-1|,|y+1|,2\}$. The exact value is $\min\{h_t(|x-1|),h_t(|x+1|),2\}$, where $h_t(d)=d^2/(2t)$ for $d\leq t$ and $d-t/2$ otherwise. At $x=0$ the two minimizing origins are $y=\pm1/2$, both of cost $3/4$. The left panel shows $g$ and its value; the right panel shows the minimized cost at $x=0$.}
{ch10-hopf-lax}
\Needspace{42mm}
\PDEAnnotation{最小化速度为何有界}{由与 $y=x$ 比较得 $L(q)-G|q|\leq L(0)$。超线性把 $q$ 限在固定球内，因此短步最小化点满足 $|x-y|=O(h)$，下接触证明中的 Taylor 余项才一致为 $o(h)$。}
